\setcounter{section}{5}

  \zwischenueberschrift{Kelengkungan Utama}

Misalkan

\mavergleichskette
{\vergleichskette
{ W
}
{ \subseteq }{ \R^n
}
{  }{
}
{    }{
}
{   }{
}
}
{}{}{}
\definitionsverweis {terbuka}{}{,}
\maabb {h} { W } { \R
} {}
suatu
\definitionsverweis {fungsi yang dua kali terdiferensialkan secara kontinu}{}{}
dan

\mavergleichskette
{\vergleichskette
{ Y
}
{ = }{ h^{-1}(c)
}
{  }{
}
{    }{
}
{   }{
}
}
{}{}{}
adalah
\definitionsverweis {serat}{}{}
di atas

\mavergleichskette
{\vergleichskette
{ c
}
{ \in }{ \R
}
{  }{
}
{    }{
}
{   }{
}
}
{}{}{,}
dengan $h$
\definitionsverweis {reguler}{}{}
pada setiap titik di $Y$. Misalkan

\mavergleichskette
{\vergleichskette
{ P
}
{ \in }{ Y
}
{  }{
}
{    }{
}
{   }{
}
}
{}{}{.}
Kita mendefinisikan berbagai konsep kelengkungan dengan mengacu pada
\definitionsverweis {pemetaan Weingarten}{}{}
\maabbdisp {L_P} {T_PY} { T_PY
} {.}

Catatan edisi: sumber memakai hipotesis $C^1$ dalam pengantar dan dua definisi berikut, meskipun pemetaan Weingarten mendiferensialkan medan normal. Ketiga hipotesis diperkuat menjadi $C^2$ agar objek yang didefinisikan benar-benar tersedia.

\inputdefinition
{}
{

Misalkan

\mavergleichskette
{\vergleichskette
{ W
}
{ \subseteq }{ \R^n
}
{  }{
}
{    }{
}
{   }{
}
}
{}{}{}
\definitionsverweis {terbuka}{}{,}
\maabb {h} { W } { \R
} {}
suatu
\definitionsverweis {fungsi yang dua kali terdiferensialkan secara kontinu}{}{}
dan

\mavergleichskette
{\vergleichskette
{ Y
}
{ = }{ h^{-1}(c)
}
{  }{
}
{    }{
}
{   }{
}
}
{}{}{}
adalah
\definitionsverweis {serat}{}{}
di atas

\mavergleichskette
{\vergleichskette
{ c
}
{ \in }{ \R
}
{  }{
}
{    }{
}
{   }{
}
}
{}{}{,}
dengan $h$
\definitionsverweis {reguler}{}{}
pada setiap titik di $Y$. Misalkan

\mavergleichskette
{\vergleichskette
{ P
}
{ \in }{ Y
}
{  }{
}
{    }{
}
{   }{
}
}
{}{}{.}
Setiap
\definitionsverweis {nilai eigen}{}{}
dari
\definitionsverweis {pemetaan Weingarten}{}{}
\maabbeledisp {L_P} {T_P Y} { T_PY
} { v } { - \left(  D_{v}  N   \right) \left(   P  \right)
} {,}
disebut suatu
\definitionswort {kelengkungan utama}{}
dari $Y$ di $P$.

}

\inputdefinition
{}
{

Misalkan

\mavergleichskette
{\vergleichskette
{ W
}
{ \subseteq }{ \R^n
}
{  }{
}
{    }{
}
{   }{
}
}
{}{}{}
\definitionsverweis {terbuka}{}{,}
\maabb {h} { W } { \R
} {}
suatu
\definitionsverweis {fungsi yang dua kali terdiferensialkan secara kontinu}{}{}
dan

\mavergleichskette
{\vergleichskette
{ Y
}
{ = }{ h^{-1}(c)
}
{  }{
}
{    }{
}
{   }{
}
}
{}{}{}
adalah
\definitionsverweis {serat}{}{}
di atas

\mavergleichskette
{\vergleichskette
{ c
}
{ \in }{ \R
}
{  }{
}
{    }{
}
{   }{
}
}
{}{}{,}
dengan $h$
\definitionsverweis {reguler}{}{}
pada setiap titik di $Y$. Misalkan

\mavergleichskette
{\vergleichskette
{ P
}
{ \in }{ Y
}
{  }{
}
{    }{
}
{   }{
}
}
{}{}{.}
Setiap
\definitionsverweis {vektor eigen}{}{}
dari
\definitionsverweis {pemetaan Weingarten}{}{}
\maabbeledisp {L_P} {T_P Y} { T_PY
} { v } { - \left(  D_{v}  N   \right) \left(   P  \right)
} {,}
disebut suatu
\definitionswort {arah kelengkungan utama}{}
dari $Y$ di $P$.

}

Istilah \stichwort {vektor kelengkungan utama} {} juga digunakan sebagai pengganti arah kelengkungan utama; istilah arah kelengkungan utama khususnya digunakan untuk vektor eigen dengan norma $1$. Berdasarkan
[[Hyperfläche/Weingartenabbildung/Selbstadjungiert/Eigenwerte/Fakt|Korolari 4.8]],
pemetaan Weingarten memiliki suatu basis ortonormal yang terdiri atas vektor-vektor eigen, yakni arah-arah kelengkungan utama. Pada umumnya, dipilih suatu basis ortonormal yang terdiri atas arah-arah kelengkungan utama. Kelengkungan-kelengkungan utama merupakan bilangan real yang sering diurutkan menurut besarnya dan ditulis sebagai

\mavergleichskettedisp
{\vergleichskette
{ \kappa_1
}
{ \leq} { \kappa_2
}
{ \leq   \ldots \leq} { \kappa_{n-1}
}
{ } {
}
{ } {
}
}
{}{}{}
dan disebut \stichwort {tupel kelengkungan utama} {.} Setiap $\kappa$ dicantumkan sebanyak dimensi ruang eigen yang bersesuaian, yaitu sebanyak
\definitionsverweis {multiplisitas geometrik}{}{nya.}

Menurut
[[Endomorphismus/Eigenwert und charakteristisches Polynom/Fakt|Teorema 23.2 (Aljabar Linear (Osnabrück 2024-2025))]],
kelengkungan-kelengkungan utama diperoleh dengan menentukan akar-akar
\definitionsverweis {polinom karakteristik}{}{}
dari
\definitionsverweis {pemetaan Weingarten}{}{,}
sedangkan
\definitionsverweis {ruang-ruang eigen}{}{}
yang bersesuaian diperoleh menggunakan
[[Lineare Abbildung/Eigenraum als Kern/Fakt|Lemma 22.1 (Aljabar Linear (Osnabrück 2024-2025))]].

\inputbeispiel{}
{

Misalkan
\mathl{a_1 , \ldots , a_{n-1}}{} bilangan-bilangan real. Kita meninjau fungsi
\maabbeledisp {f} {\R^{n-1}} {\R
} { \left(  x_1   , \,   \ldots  , \,   x_{n-1}                 \right) } { a_1x_1^2  + \cdots + a_{n-1} x_{n-1}^2
} {.}
Menurut
[[Differenzierbare Funktion/Graph/Hyperfläche/Weingartenabbildung/Fakt|Lemma 4.9]],
\definitionsverweis {pemetaan Weingarten}{}{}
dari
\definitionsverweis {grafik}{}{}
$f$ di titik nol dijelaskan oleh
\definitionsverweis {matriks Hesse}{}{}
dari $f$, yaitu oleh

\mathdisp {\begin{pmatrix}  2a_1  & 0 & 0 & 0 \\ 0 &  2a_2  & 0  & 0 \\ 0 & 0 &   \ddots & 0 \\ 0 & 0 & 0 &  2a_{n-1}  \end{pmatrix}} { . }

Jadi,
\definitionsverweis {kelengkungan-kelengkungan utama}{}{}
grafik tersebut di titik nol adalah $2a_1  , \ldots , 2a_{n-1}$ dan
\definitionsverweis {arah-arah kelengkungan utama}{}{}
diberikan oleh vektor-vektor standar. Secara khusus, setiap tupel bilangan real muncul sebagai tupel kelengkungan utama suatu hipermuka terdiferensialkan.

}

  \zwischenueberschrift{Kelengkungan pada Permukaan}

Sekarang kita menelaah lebih saksama konsep-konsep kelengkungan pada permukaan terdiferensialkan di ruang.

\inputdefinition
{}
{

Misalkan

\mavergleichskette
{\vergleichskette
{ W
}
{ \subseteq }{ \R^3
}
{  }{
}
{    }{
}
{   }{
}
}
{}{}{}
\definitionsverweis {terbuka}{}{,}
\maabb {h} { W } { \R
} {}
suatu
\definitionsverweis {fungsi yang dua kali terdiferensialkan secara kontinu}{}{}
dan

\mavergleichskette
{\vergleichskette
{ Y
}
{ = }{ h^{-1}(c)
}
{  }{
}
{    }{
}
{   }{
}
}
{}{}{}
adalah serat di atas

\mavergleichskette
{\vergleichskette
{ c
}
{ \in }{ \R
}
{  }{
}
{    }{
}
{   }{
}
}
{}{}{,}
dengan $h$
\definitionsverweis {reguler}{}{}
pada setiap titik di $Y$, dan misalkan

\mavergleichskette
{\vergleichskette
{ P
}
{ \in }{ Y
}
{  }{
}
{    }{
}
{   }{
}
}
{}{}{.}
\definitionsverweis {Rata-rata aritmetika}{}{}

\mathl{{ \frac{   \kappa_1+\kappa_2  }{  2 } }}{} dari kedua
\definitionsverweis {kelengkungan utama}{}{}
di $P$ disebut
\definitionswort {kelengkungan rata-rata}{}
permukaan $Y$ di $P$.

}

\inputdefinition
{}
{

Misalkan

\mavergleichskette
{\vergleichskette
{ W
}
{ \subseteq }{ \R^3
}
{  }{
}
{    }{
}
{   }{
}
}
{}{}{}
\definitionsverweis {terbuka}{}{,}
\maabb {h} { W } { \R
} {}
suatu
\definitionsverweis {fungsi yang dua kali terdiferensialkan secara kontinu}{}{}
dan

\mavergleichskette
{\vergleichskette
{ Y
}
{ = }{ h^{-1}(c)
}
{  }{
}
{    }{
}
{   }{
}
}
{}{}{}
adalah serat di atas

\mavergleichskette
{\vergleichskette
{ c
}
{ \in }{ \R
}
{  }{
}
{    }{
}
{   }{
}
}
{}{}{,}
dengan $h$
\definitionsverweis {reguler}{}{}
pada setiap titik di $Y$, dan misalkan

\mavergleichskette
{\vergleichskette
{ P
}
{ \in }{ Y
}
{  }{
}
{    }{
}
{   }{
}
}
{}{}{.}
Hasil kali
\mathl{\kappa_1 \cdot \kappa_2}{} dari kedua
\definitionsverweis {kelengkungan utama}{}{}
di $P$ disebut
\definitionswort {kelengkungan Gauss}{}
permukaan $Y$ di $P$.

}

Kelengkungan Gauss adalah
\definitionsverweis {determinan}{}{}
dari pemetaan Weingarten.

\inputbeispiel{}
{

Pada permukaan bola berjari-jari $r$,
\definitionsverweis {pemetaan Weingarten}{}{}
\zusatzklammer {terhadap medan normal satuan yang mengarah ke dalam} {} {}
di setiap titik, menurut
[[Kugel/Radius/Weingartenabbildung/Beispiel|Contoh 4.4]],
adalah homoteti skalar dengan faktor ${ \frac{   1  }{ r } }$. Oleh karena itu, ${ \frac{   1  }{ r } }$ merupakan satu-satunya
\definitionsverweis {kelengkungan utama}{}{}
dengan multiplisitas $2$, dan
\definitionsverweis {kelengkungan Gauss}{}{}
adalah
\mathl{{ \frac{   1  }{ r^2 } }}{.}

}

\inputbeispiel{}
{

Untuk hiperboloid satu lembar

\mavergleichskettedisp
{\vergleichskette
{ Y
}
{ =} { \left\{ (x,y,z)  \mid   x^2+y^2-z^2  = 1         \right\}
}
{ \subseteq} { \R^3
}
{ } {
}
{ } {
}
}
{}{}{}
berdasarkan perhitungan dalam
[[Einschaliges Hyperboloid/Weingartenabbildung/Beispiel|Contoh 4.5]],
pada suatu titik

\mavergleichskette
{\vergleichskette
{ (x,y,z)
}
{ \in }{ Y
}
{  }{
}
{    }{
}
{   }{
}
}
{}{}{}
kedua
\definitionsverweis {kelengkungan utama}{}{}
adalah
\mathkor {} {- \left(  x^2+y^2+z^2  \right)^{- { \frac{   1  }{  2 } } }} {dan} {\left(  x^2+y^2+z^2  \right)^{- { \frac{   3  }{  2 } } }} {,}
sehingga
\definitionsverweis {kelengkungan Gauss}{}{}
adalah

\mavergleichskettedisphandlinks
{\vergleichskettedisphandlinks
{ - \left(  x^2+y^2+z^2  \right)^{- { \frac{   1  }{  2 } } }  \cdot \left(  x^2+y^2+z^2  \right)^{- { \frac{   3  }{  2 } } }
}
{ =} { - \left(  x^2+y^2+z^2  \right)^{- 2 }
}
{ =} { -  { \frac{   1  }{  \left(  x^2+y^2+z^2  \right)^{ 2 }   } }
}
{ } {
}
{ } {
}
}
{}{}{.}

}

\inputbeispiel{}
{

Misalkan $Y$ adalah
\definitionsverweis {grafik}{}{}
suatu fungsi yang dua kali terdiferensialkan secara kontinu
\maabbele {f} { V } { \R
} {(x,y)} { f(x,y)
} {,}
pada suatu himpunan terbuka

\mavergleichskette
{\vergleichskette
{ V
}
{ \subseteq }{ \R^2
}
{  }{
}
{    }{
}
{   }{
}
}
{}{}{.}
Menurut
[[Differenzierbare Funktion/Graph/Hyperfläche/Weingartenabbildung/Fakt|Lemma 4.9]],
\definitionsverweis {pemetaan Weingarten}{}{}
\zusatzklammer {terhadap medan normal satuan yang mengarah ke atas} {} {}
di setiap titik

\mavergleichskette
{\vergleichskette
{ P
}
{ = }{ (Q,f(Q))
}
{  }{
}
{    }{
}
{   }{
}
}
{}{}{}
memiliki bentuk berikut. Tuliskan
\mathl{g=\operatorname{Grad} f(Q)}{,}
\mathl{G=I_2+gg^{\mathsf T}}{,}
\mathl{H=\operatorname{Hess} f(Q)}{,}
dan
\mathl{\omega=\sqrt{1+\lVert g\rVert^2}}{.}
Dalam basis koordinat grafik, $G$ adalah matriks metrik dan matriks $A$ dari pemetaan Weingarten adalah

\mathdisp {A=G^{-1}{ \frac{ 1 }{ \omega } }H} { . }

Jadi, koordinat $u$ dari suatu arah kelengkungan utama memenuhi masalah nilai eigen umum

\mathdisp {Hu=\kappa\,\omega Gu} { . }

Hanya apabila $g=0$ masalah ini mereduksi menjadi masalah nilai eigen biasa untuk matriks Hesse. Untuk permukaan ini,
\definitionsverweis {kelengkungan Gauss}{}{}
adalah

\mathdisp {K={ \frac{ \det H }{ \omega^4 } }={ \frac{ \det\!\left(\operatorname{Hess}f(Q)\right) }{ \left(1+\lVert\operatorname{Grad}f(Q)\rVert^2\right)^2 } }} { . }

Catatan edisi: sumber menunjuk contoh bola yang tidak terkait dan menghilangkan faktor invers metrik $G^{-1}$. Akibatnya, pernyataan sumber tentang arah utama dan penyebut kelengkungan Gauss tidak benar secara umum; pernyataan itu tereduksi ke rumus yang benar pada titik kritis $f$. Rumus di atas menerapkan Lemma 4.9 yang telah diperbaiki pada Unit 4.

}

\inputbeispiel{}
{

Misalkan $I$ suatu
\definitionsverweis {interval terbuka}{}{}
dan
\maabb {f} { I } { \R_+
} {}
suatu
\definitionsverweis {fungsi yang dua kali terdiferensialkan secara kontinu}{}{,}
misalkan $Y$ adalah
\definitionsverweis {permukaan putar}{}{}
yang bersesuaian
\zusatzklammer {mengelilingi sumbu $x$} {} {}
dengan
\definitionsverweis {grafik}{}{}
$f$, dan misalkan

\mavergleichskettedisp
{\vergleichskette
{ P
}
{ =} { \begin{pmatrix}  x_0  \\f(x_0)\\  0   \end{pmatrix}
}
{ \in} { Y
}
{ } {
}
{ } {
}
}
{}{}{}
\zusatzklammer {yang bukan merupakan pembatasan berarti} {} {.}
Kita meninjau bagian permukaan putar yang, pada suatu lingkungan $(x_0,0)$ yang termuat dalam
\mathl{\left\{ (x,z)   \mid   z^2 < f(x)^2        \right\}=\left\{(x,z)\mid |z|<f(x)\right\}}{,} dapat ditulis sebagai grafik dari

\mavergleichskettedisp
{\vergleichskette
{ g(x,z)
}
{ =} { \sqrt{ f(x)^2-z^2 }
}
{ =} { \left(  f(x)^2-z^2   \right)^{1/2}
}
{ } {
}
{ } {
}
}
{}{}{.}
Matriks Jacobi dari $g$ adalah

\mathdisp {\left(  \left(  f(x)^2-z^2  \right)^{-1/2}  f(x) f'(x)   , \,   - z \left(  f(x)^2-z^2  \right)^{-1/2}                   \right)} {  }

Untuk meringkas tampilan matriks berikut, tulis $s=f(x)^2-z^2>0$. Matriks Hesse dari $g$ adalah

\mathdisp {\begin{pmatrix}  -s^{-3/2} f(x)^2 f'(x)^2 + s^{-1/2} \left(  f'(x)^2 + f(x) f^{\prime \prime} (x)   \right)   &   z s^{-3/2} f(x) f'(x)   \\  z s^{-3/2} f(x) f'(x)  &  -s^{-1/2}  -z^2 s^{-3/2}  \end{pmatrix}} {  }

atau merupakan hasil kali skalar
\mathl{\left(  f(x)^2-z^2  \right)^{-3/2}}{} dengan matriks

\mavergleichskettealigndrucklinks
{\vergleichskettealigndrucklinks
{ \begin{pmatrix}  -  f(x)^2 f'(x)^2 + \left(  f(x)^2-z^2  \right) \left(  f'(x)^2 + f(x) f^{\prime \prime} (x)   \right)   &   z  f(x) f'(x)   \\  z f(x) f'(x)  &  - \left(  f(x)^2-z^2  \right)  -z^2   \end{pmatrix}
}
{ =} { \begin{pmatrix}  f(x)^3 f^{\prime \prime} (x)  -z^2 f'(x)^2    -z^2 f(x) f^{\prime \prime} (x)    &   z  f(x) f'(x)   \\  z f(x) f'(x)  &  - f(x)^2  \end{pmatrix}
}
{ } {
}
{ } {
}
{ } {
}
}
{}{}{.}
Untuk titik $P$ yang dipilih, hasil ini adalah

\mavergleichskettedisp
{\vergleichskette
{ { \frac{   1  }{  f(x_0)^3 } } \begin{pmatrix}  f(x_0)^3 f^{\prime \prime} (x_0)   &   0   \\  0  &  - f(x_0)^2  \end{pmatrix}
}
{ =} { \begin{pmatrix}  f^{\prime \prime} (x_0)   &   0   \\  0  &  -  { \frac{   1  }{ f(x_0) } }   \end{pmatrix}
}
{ } {
}
{ } {
}
{ } {
}
}
{}{}{.}
Menurut
[[Differenzierbare Funktion/Graph/Hyperfläche/Weingartenabbildung/Fakt|Lemma 4.9]],
dengan
\mathl{\omega=\sqrt{1+f'(x_0)^2}}{,}
faktor invers metrik grafik juga harus diperhitungkan. Pemetaan Weingarten di $P$ adalah

\mathdisp {\begin{pmatrix}  { \frac{ f^{\prime \prime}(x_0) }{ \omega^3 } } & 0 \\ 0 & -{ \frac{ 1 }{ f(x_0)\omega } } \end{pmatrix}} { . }

Dengan demikian, kedua
\definitionsverweis {kelengkungan utama}{}{}
adalah
\mathkor {} {{ \frac{   f^{\prime \prime} (x_0)  }{  \left(1+f'(x_0)^2\right)^{3/2}  } }} {dan} {- { \frac{   1  }{  f(x_0) \sqrt{ 1 +f'(x_0)^2 }  } }} {,}
sedangkan arah-arah kelengkungan utama adalah vektor-vektor standar atau citranya
\mathkor {} {\begin{pmatrix}  1  \\f'(x_0)\\  0   \end{pmatrix}} {dan} {\begin{pmatrix}  0  \\ 0\\  1   \end{pmatrix}} {}
di ruang singgung $T_PY$. Jadi, arah-arah kelengkungan utama membentang sepanjang grafik $f$ dan sepanjang gerak melingkar rotasi.
\definitionsverweis {Kelengkungan Gauss}{}{}
adalah

\mathdisp {{ \frac{   - f^{\prime \prime} (x_0)  }{  f(x_0)  \left(  1 +f'(x_0)^2  \right)^2    } }} { . }

Catatan edisi: sumber hanya mengasumsikan keterdiferensialan meskipun memakai $f''$, menuliskan daerah grafik yang keliru, dan kembali menghilangkan faktor invers metrik. Hipotesis, daerah, dan tiga rumus kelengkungan di atas telah diperbaiki.

}

  \zwischenueberschrift{Kelengkungan Normal}

\inputdefinition
{}
{

Misalkan

\mavergleichskette
{\vergleichskette
{ W
}
{ \subseteq }{ \R^n
}
{  }{
}
{    }{
}
{   }{
}
}
{}{}{}
\definitionsverweis {terbuka}{}{,}
\maabb {h} { W } { \R
} {}
suatu
\definitionsverweis {fungsi yang dua kali terdiferensialkan secara kontinu}{}{}
dan

\mavergleichskette
{\vergleichskette
{ Y
}
{ = }{ h^{-1}(c)
}
{  }{
}
{    }{
}
{   }{
}
}
{}{}{}
adalah
\definitionsverweis {serat}{}{}
di atas

\mavergleichskette
{\vergleichskette
{ c
}
{ \in }{ \R
}
{  }{
}
{    }{
}
{   }{
}
}
{}{}{,}
dengan $h$
\definitionsverweis {reguler}{}{}
pada setiap titik di $Y$. Misalkan

\mavergleichskette
{\vergleichskette
{ P
}
{ \in }{ Y
}
{  }{
}
{    }{
}
{   }{
}
}
{}{}{.}
Untuk suatu vektor singgung ternormalisasi

\mavergleichskette
{\vergleichskette
{ v
}
{ \in }{ T_PY
}
{  }{
}
{    }{
}
{   }{
}
}
{}{}{}
nilai
\mathl{\left\langle  v  ,  L_P(v)  \right\rangle}{} disebut
\definitionswort {kelengkungan normal}{}
dari $Y$ di $P$ dalam arah $v$. Nilai tersebut dilambangkan dengan

\mavergleichskette
{\vergleichskette
{ \kappa(P,v)
}
{ = }{ \kappa(v)
}
{  }{
}
{    }{
}
{   }{
}
}
{}{}{}
.

}

Menurut
[[Hyperfläche/Weingartenabbildung/Zweite Ableitung/Skalarprodukt/Fakt|Lemma 4.6]],
kelengkungan normal sama dengan

\mavergleichskettedisp
{\vergleichskette
{ \left\langle  v  ,  L_P(v)  \right\rangle
}
{ =} { \left\langle  \gamma^{\prime \prime} (0)  ,  N(P)  \right\rangle
}
{ } {
}
{ } {
}
{ } {
}
}
{}{}{,}
dengan $\gamma$ suatu realisasi kurva dari vektor singgung $v$. Jadi, kelengkungan normal mengukur komponen normal percepatan kurva tersebut.

__NOEDITSECTION__

\inputfaktbeweis
{Hyperfläche/Normalkrümmung/Berechnung/Fakt}
{Lemma}
{}
{

\faktsituation {Misalkan

\mavergleichskette
{\vergleichskette
{ W
}
{ \subseteq }{ \R^n
}
{  }{
}
{    }{
}
{   }{
}
}
{}{}{}
\definitionsverweis {terbuka}{}{,}
\maabb {h} { W } { \R
} {}
suatu
\definitionsverweis {fungsi yang dua kali terdiferensialkan secara kontinu}{}{}
dan

\mavergleichskette
{\vergleichskette
{ Y
}
{ = }{ h^{-1}(c)
}
{  }{
}
{    }{
}
{   }{
}
}
{}{}{}
adalah
\definitionsverweis {serat}{}{}
di atas

\mavergleichskette
{\vergleichskette
{ c
}
{ \in }{ \R
}
{  }{
}
{    }{
}
{   }{
}
}
{}{}{,}
dengan $h$
\definitionsverweis {reguler}{}{}
pada setiap titik di $Y$. Misalkan

\mavergleichskette
{\vergleichskette
{ P
}
{ \in }{ Y
}
{  }{
}
{    }{
}
{   }{
}
}
{}{}{.}
Misalkan
\mathl{\kappa_1 , \ldots , \kappa_{n-1}}{} adalah
\definitionsverweis {kelengkungan-kelengkungan utama}{}{}
dari $Y$ di $P$ dan
\mathl{v_1 , \ldots , v_{n-1}}{} adalah
\definitionsverweis {arah-arah kelengkungan utama}{}{}
\definitionsverweis {ternormalisasi}{}{ yang bersesuaian dan membentuk suatu basis ortonormal.}}
\faktfolgerung {Maka
\definitionsverweis {kelengkungan normal}{}{}
dari $Y$ di $P$ dalam arah

\mavergleichskette
{\vergleichskette
{ v
}
{ \in }{ T_PY
}
{  }{
}
{    }{
}
{   }{
}
}
{}{}{}
yang memenuhi $\lVert v\rVert=1$ adalah

\mavergleichskettedisp
{\vergleichskette
{ \kappa(v)
}
{ =} { \sum_{i = 1}^{n-1} \kappa_i \left\langle  v  , v_i   \right\rangle^2
}
{ } {
}
{ } {
}
{ } {
}
}
{}{}{.}}
Untuk vektor tak nol yang belum dinormalisasi, bentuk kuadratik pada ruas kanan memenuhi

\mathdisp {\sum_{i=1}^{n-1}\kappa_i\left\langle v,v_i\right\rangle^2=\lVert v\rVert^2\kappa\!\left({ \frac{v}{\lVert v\rVert} }\right)} { . }

Catatan edisi: sumber hanya menyebut vektor-vektor utama ``ternormalisasi'', padahal bukti memakai ortogonalitas, dan tidak mensyaratkan arah $v$ berpanjang satu. Kedua hipotesis itu dibuat eksplisit di sini.
\faktzusatz {}
\faktzusatz {}

}
{

Dengan

\mavergleichskette
{\vergleichskette
{ v
}
{ = }{ \sum_{i = 1}^{n-1} a_i v_i
}
{  }{
}
{    }{
}
{   }{
}
}
{}{}{}
dengan menggunakan ortogonalitas, diperoleh

\mavergleichskettealign
{\vergleichskettealign
{ \kappa(v)
}
{ =} { \left\langle  v  ,  L_P(v)  \right\rangle
}
{ =} { \left\langle  \sum_{i = 1}^{n-1} a_i v_i  ,  L_P \left(  \sum_{i = 1}^{n-1} a_i v_i  \right)   \right\rangle
}
{ =} { \left\langle  \sum_{i = 1}^{n-1} a_i v_i  ,  \sum_{i = 1}^{n-1} \kappa_i a_i v_i   \right\rangle
}
{ =} { \sum_{i = 1}^{n-1} \kappa_i a_i^2
}
}
{
\vergleichskettefortsetzungalign
{ =} { \sum_{i = 1}^{n-1} \kappa_i \left\langle  v  , v_i   \right\rangle^2
}
{ } {}
{ } {}
{ } {}
}
{}{.}

}

\bild{ \begin{center}
\includegraphics[width=5.5cm]{ \bildeinlesungsvg {Minimal surface curvature planes-de} {svg} }
\end{center}
\bildtext {Bidang-bidang kelengkungan utama, vektor normal, dan bidang singgung pada sebuah permukaan minimal. Label dalam berkas sumber dijelaskan seluruhnya oleh keterangan ini.} }

\bildlizenz { Minimal surface curvature planes-de.svg } {} {Eric Gaba} {Commons} {CC-by-sa 3.0} {}

\inputdefinition
{}
{

Misalkan

\mavergleichskette
{\vergleichskette
{ W
}
{ \subseteq }{ \R^n
}
{  }{
}
{    }{
}
{   }{
}
}
{}{}{}
\definitionsverweis {terbuka}{}{,}
\maabb {h} { W } { \R
} {}
suatu
\definitionsverweis {fungsi yang dua kali terdiferensialkan secara kontinu}{}{}
dan

\mavergleichskette
{\vergleichskette
{ Y
}
{ = }{ h^{-1}(c)
}
{  }{
}
{    }{
}
{   }{
}
}
{}{}{}
adalah
\definitionsverweis {serat}{}{}
di atas

\mavergleichskette
{\vergleichskette
{ c
}
{ \in }{ \R
}
{  }{
}
{    }{
}
{   }{
}
}
{}{}{,}
dengan $h$
\definitionsverweis {reguler}{}{}
pada setiap titik di $Y$. Misalkan

\mavergleichskette
{\vergleichskette
{ P
}
{ \in }{ Y
}
{  }{
}
{    }{
}
{   }{
}
}
{}{}{.}
Misalkan

\mavergleichskette
{\vergleichskette
{ v
}
{ \in }{ T_PY
}
{  }{
}
{    }{
}
{   }{
}
}
{}{}{}
suatu
\definitionsverweis {vektor singgung}{}{}
yang berbeda dari $0$, dan misalkan

\mavergleichskette
{\vergleichskette
{ u
}
{ \in }{ N_PY
}
{  }{
}
{    }{
}
{   }{
}
}
{}{}{}
suatu
\definitionsverweis {vektor normal}{}{.}
Vektor tersebut berbeda dari $0$. Bidang $\R u + \R v$ disebut
\definitionswort {bidang normal}{}
dari $Y$ melalui $P$.

}

Setiap vektor singgung

\mavergleichskette
{\vergleichskette
{ v
}
{ \neq }{ 0
}
{  }{
}
{    }{
}
{   }{
}
}
{}{}{}
menentukan sebuah bidang normal $E$ secara tunggal, yang biasanya dipandang sebagai
\mathl{P+E}{}.
__NOEDITSECTION__

\inputfaktbeweis
{Hyperfläche/Normalebene/Kurvenschnitt/Regularität/Fakt}
{Lemma}
{}
{

\faktsituation {Misalkan

\mavergleichskette
{\vergleichskette
{ W
}
{ \subseteq }{ \R^n
}
{  }{
}
{    }{
}
{   }{
}
}
{}{}{}
\definitionsverweis {terbuka}{}{,}
\maabb {h} { W } { \R
} {}
suatu
\definitionsverweis {fungsi yang dua kali terdiferensialkan secara kontinu}{}{}
dan

\mavergleichskette
{\vergleichskette
{ Y
}
{ = }{ h^{-1}(c)
}
{  }{
}
{    }{
}
{   }{
}
}
{}{}{}
adalah
\definitionsverweis {serat}{}{}
di atas

\mavergleichskette
{\vergleichskette
{ c
}
{ \in }{ \R
}
{  }{
}
{    }{
}
{   }{
}
}
{}{}{,}
dengan $h$
\definitionsverweis {reguler}{}{}
pada setiap titik di $Y$. Misalkan

\mavergleichskette
{\vergleichskette
{ P
}
{ \in }{ Y
}
{  }{
}
{    }{
}
{   }{
}
}
{}{}{.}
Misalkan

\mavergleichskette
{\vergleichskette
{ E
}
{ \subseteq }{ \R^n
}
{  }{
}
{    }{
}
{   }{
}
}
{}{}{}
suatu
\definitionsverweis {bidang normal}{}{}
melalui $P$ pada $Y$.}
\faktfolgerung {Maka irisan
\mathl{Y \cap (P+E)}{} merupakan suatu kurva bidang di $P+E$ yang reguler pada titik $P$.}
\faktzusatz {}
\faktzusatz {}

}
{

Setelah melakukan translasi, kita dapat mengasumsikan bahwa $P$ adalah titik nol ruang tersebut. Diferensial total
\maabbdisp {\left(D h \right)_{ 0}} {\R^n } {\R
} {}
bersifat surjektif dan kernelnya adalah ruang singgung $T_PY$. Bidang $E$ memuat vektor-vektor normal $\neq 0$ yang tidak berada dalam kernel. Oleh karena itu, diferensial total dari pemetaan komposit

\mathdisp {E \supseteq E \cap W \longrightarrow W \stackrel{h}{\longrightarrow} \R} {  }

bersifat surjektif di titik $0$. Dengan demikian, teorema fungsi implisit dapat diterapkan, dan diperoleh bahwa serat tersebut merupakan kurva yang reguler di titik $0$.

}

Perhatikan bahwa irisan kurva tersebut tidak harus reguler di setiap titik.
__NOEDITSECTION__

\inputfaktbeweis
{Hyperfläche/Normalkrümmung/Normalebene/Kurvenschnitt/Fakt}
{Satz}
{}
{

\faktsituation {Misalkan

\mavergleichskette
{\vergleichskette
{ W
}
{ \subseteq }{ \R^n
}
{  }{
}
{    }{
}
{   }{
}
}
{}{}{}
\definitionsverweis {terbuka}{}{,}
\maabb {h} { W } { \R
} {}
suatu
\definitionsverweis {fungsi yang dua kali terdiferensialkan secara kontinu}{}{}
dan

\mavergleichskette
{\vergleichskette
{ Y
}
{ = }{ h^{-1}(c)
}
{  }{
}
{    }{
}
{   }{
}
}
{}{}{}
adalah
\definitionsverweis {serat}{}{}
di atas

\mavergleichskette
{\vergleichskette
{ c
}
{ \in }{ \R
}
{  }{
}
{    }{
}
{   }{
}
}
{}{}{,}
dengan $h$
\definitionsverweis {reguler}{}{}
pada setiap titik di $Y$. Misalkan

\mavergleichskette
{\vergleichskette
{ P
}
{ \in }{ Y
}
{  }{
}
{    }{
}
{   }{
}
}
{}{}{.}
Misalkan

\mavergleichskette
{\vergleichskette
{ E
}
{ \subseteq }{ \R^n
}
{  }{
}
{    }{
}
{   }{
}
}
{}{}{}
suatu
\definitionsverweis {bidang normal}{}{}
melalui $P$ pada $Y$. Ambil suatu vektor satuan
\mathl{v\in E\cap T_PY}{,}
dan orientasikan bidang $P+E$ dengan menetapkan pasangan $(v,N(P))$ sebagai basis berorientasi positif. Dengan demikian, normal satuan kurva irisan di dalam bidang yang bersesuaian dengan arah singgung $v$ adalah $N(P)$.}
\faktfolgerung {Maka
\definitionsverweis {kelengkungan normal}{}{}
dari $Y$ di $P$ dalam arah $v$ sama dengan kelengkungan bertanda
\definitionsverweis {kelengkungan}{}{}
kurva bidang

\mavergleichskette
{\vergleichskette
{ Y \cap (P+E)
}
{ \subseteq }{ P+E
}
{ \cong }{ \R^2
}
{    }{
}
{   }{
}
}
{}{}{}
di titik $P$. Tanpa pilihan orientasi yang kompatibel, hanya nilai mutlak kedua kelengkungan yang ditentukan secara invarian.}
\faktzusatz {}
\faktzusatz {}

}
{

Untuk menyederhanakan notasi, kita mentranslasikan $P$ ke titik nol. Menurut
[[Hyperfläche/Normalebene/Kurvenschnitt/Regularität/Fakt|Lemma 5.13]],
irisan $Y \cap E$ merupakan kurva yang reguler di $0$. Misalkan
\maabbdisp {\gamma} { I } { Y \cap E
} {}
suatu realisasi yang dua kali terdiferensialkan dan berparametrisasi panjang busur dengan
\mathl{\gamma(0)=0}{,}
\mathl{\gamma'(0)=v}{,}
dan normal satuan kurva $\gamma$ di dalam bidang yang kompatibel sama dengan $N(0)$. Menurut
[[Hyperfläche/Weingartenabbildung/Zweite Ableitung/Skalarprodukt/Fakt|Lemma 4.6]],
kelengkungan normal sama dengan
\mathl{\left\langle  \gamma^{\prime \prime}(0) , N(0)  \right\rangle}{.} Karena bidang normal memuat vektor normal satuan $N(0)$, seluruh perhitungan berlangsung di bidang $E$. Jadi, pernyataan tersebut mengikuti dari
[[Ebene reguläre Kurve/Krümmung/Skalarprodukt/Fakt|Lemma 3.8]].

Catatan edisi: sumber tidak menetapkan orientasi irisan bidang, sehingga kesamaan kelengkungan bertanda bersifat ambigu. Konvensi orientasi di atas menghilangkan ambiguitas tersebut.

}

 [[Kategorie:Latexseite]]
